\section{Ontotheology}

Before moving on to the discussion of the One and the Many, it is necessary to clarify some terms. The position of Gornahoor is ontotheological\footnote{\url{https://gornahoor.net/?p=1299}}, which we believe is consistent with Guenonian metaphysics and the Western tradition. By this we mean the deep relationship of God or Absolute, Being, and Logos. However, unlike Western philosophy (such as Aristotle, Aquinas, Hegel), which is incomplete since it equates the Absolute with Being, metaphysics does not stop at Being since that doesn't take into account the Infinitude of Total Possibility. However, this is not our concern at the moment.

Ontotheology met with criticism from Kant, who regarded it as abstract, to Heidegger, who considered it the ``forgetfulness of Being''. It is abstract if the only experience admitted is synthetic, but we recognize a form of knowledge beyond abstract thought, that is, a direct intuition of higher states.

If all experience is synthetic, then we are faced with the representation of the World as a series of discrete experiences. From those experiences, the philosopher then constructs a mental abstraction. In contrast, the metaphysician understands experience as a ``seamless, indefinable, alogical continuum'' (Woodroffe), so it is grasped as a whole, beyond time and space.

\paragraph{Sat, Chit, Ananda}


In Tantra, the Universe of Experience has three aspects common to all objective experience: Being (the object is), Consciousness (it is known), and Pleasure (the experience is pleasant). Hence, Being is inseparable from the knowing subject, as ``to be'' is ``to be conscious''. A higher level of being means a deeper level of consciousness. The ability to realize higher states is a measure of Power. Higher states involve more wholeness and fullness, or limitless which is experienced as bliss. The diminution of being --- limitation or privation --- is experienced as pain.

\paragraph{Existence}

Existence is the manifestation of possibilities and a World is the manifestation of compossibilities. A World is one state of Total Existence, which comprises all possibilities of manifestation, since the possible must become real in some way. However, Being is the principle of manifestation, so Being never is manifested as such. Thus, from the metaphysical point of view, God does not ``exist'' in this sense. In passing, we can note that both Guenon and Woodroffe include the idea of multiple worlds as Traditional teachings; this is the metaphysical analog of the ``multiverse''.

\paragraph{Logos}


The Primordial Tradition includes the existence of a ``cosmic and juridical order''\footnote{\url{https://gornahoor.net/?p=2053}}, the two aspects of which correspond to the Logos and the implementation of justice in the world. This is co-extensive with Being, otherwise it would be a limitation on Being. It is the principle of compossibility, as only the rational is real, and what is possible is related to everything else in manifestation.

\paragraph{Actualization}


Increase in Being is far from abstract. Being is not known without a corresponding change in state. Thus, Being is known most fully when it is actualized in the world through power. As an example, let us take a widely known tune, say, \emph{Somewhere over the Rainbow}. As a phrase, it is an abstract idea. You may ``hear'' the song in your head, giving it some degree of existence or perhaps hum it, giving it a little more. A band will perform the song, giving it full existence. Note that as real, the song is manifested over time, although it is a single idea. It cannot happen ``all at once'' or else it would be a muddle of cacophony.

\paragraph{References}

\emph{The World as Power}, John Woodroffe

\emph{The Multiple States of Being}, Rene Guenon

\emph{Hegel and the Hermetic Tradition}, Glenn Magee


\flrightit{Posted on 2011-04-06 by Cologero}

